\section{Theoretical Assumptions}
\label{sec:theoretical_assumptions}

The framework relies on several assumptions that define the intended interpretation of the intrinsic reward.
These assumptions do not guarantee global optimality of the learned policy, but they specify the conditions under which feature-space impact, learning progress, and gating are expected to provide useful exploration guidance.

First, the environment is assumed to provide bounded episodic interaction.
Extrinsic rewards, intrinsic rewards after clipping, and episode horizons are finite, so the discounted returns used by the policy optimizer are well defined.
This assumption is necessary because the augmented reward in Equation~\ref{eq:reward_augmentation} changes the value targets and advantages used during training.

Second, the learned representation is assumed to preserve action-relevant transition structure.
The latent vector $z_t$ need not reconstruct the observation or represent the full environment state, but distances in latent space should be meaningful enough for $\lVert z_{t+1}-z_t\rVert_2$ to indicate transition impact.
The inverse-dynamics objective supports this assumption by encouraging features that retain information about actions that caused observed transitions.

Third, the forward prediction error is assumed to be a usable proxy for local model competence.
A decreasing forward error in a region is interpreted as evidence that the model is learning the local dynamics.
The framework does not assume that all uncertainty is reducible.
Instead, it explicitly allows regions with high error and low progress, which are candidates for gate suppression.

Fourth, region-local statistics are assumed to receive enough samples to become meaningful.
The online partition should be neither so coarse that unrelated dynamics are mixed nor so fine that most regions lack sufficient data.
The capacity and maximum-depth parameters of the partition control this tradeoff.
The requirement of a minimum number of visited regions before gating further prevents early median estimates from controlling exploration prematurely.

Fifth, the exponential moving averages are assumed to operate on a slower time scale than individual transition noise.
The inequality $\beta_{\mathrm{short}}<\beta_{\mathrm{long}}$ ensures that the short-horizon estimate reacts faster than the long-horizon estimate.
Learning progress is therefore interpreted as a relative improvement signal, not as an absolute measure of environment difficulty.

Sixth, component-wise normalization is assumed to stabilize relative scaling without changing the qualitative meaning of the components.
The normalized quantities $\widetilde{I}_t$ and $\widetilde{\mathrm{LP}}_t$ are scale-adjusted versions of impact and learning progress.
The weights $\alpha_{\mathrm{impact}}$ and $\alpha_{\mathrm{LP}}$ then determine the intended tradeoff rather than compensating for arbitrary units.

Finally, the framework assumes that intrinsic reward is a training aid rather than a separate task objective.
The augmented objective may alter the policy optimization path and is not claimed to preserve the optimal policy of the original MDP in the strict reward-shaping sense.
The method is therefore evaluated by whether the resulting policy achieves higher or more reliable extrinsic return, not by the magnitude of accumulated intrinsic reward.